\scanpage{32}
\mspage{9}
\mspage{31}

est essentiellement constant, \deleted{pour} à fortiori satisfait la condition de Mittag-Leffler, donc (compte tenu que $\widehat F=\lim F_m$) on conclut (ii) par EGA $0_{\mathrm{III}}$ 13/3. D'autre part (i) en résulte trivialement, compte tenu de 2.3.

\uline{Corollaire} 2.4. Soient $f:X\quad S$ un morphisme projectif et plat, avec $S$ localement noethérien, $\underline O_X(1)$ un Module inversible sur $X$, ample rela\deleted{\unsure{ti}}tivement à $S$, $t$ une section de ce module, qui soit $\underline O_X$-régulière, $Y$ \deleted{un} le sous-\unsure{présbhéma|préschéma} des zéros de $t$, $\widehat X$ le complété formel de $X$ le long de $Y$. On suppose que pour tout $s\quad S$, $X_s$ est de profondeur $\quad1$ (resp. de profondeur $\quad2$) en ses points fermés. Alors pour tout voisinage ouvert $U$ de $Y$, le foncteur
\[
F\qquad\widehat F
\]
de la catégorie des Modules cohérents localement libres sur $U$, dans la catégorie des Modules cohérents localement libres sur $\widehat X$, est fidèle (resp. pleinement fidèle, i.e. la condition \deleted{(Lef)} de Lefschetz (Lef) de l'exposé X N°2 est vérifiée).

Introduisons \deleted{le Module} pour deux Modules localement libres $F$ et $G$ sur $U$ le Module
\[
H=\underline{\operatorname{Hom}}_{\underline O_U}(F,G),
\]
on est ramené à prouver que l'homomorphisme canonique
\[
H^0(U,H)\qquad H^0(\widehat X,\widehat H)
\]
est injectif (resp. bijectif). Or les Modules $H_t$ sont de profondeur $\quad1$ (resp. $\quad2$) aux points fermés de $X_y$, on peut donc appliquer 2.1., qui implique la conclusion 2.4. dans le cas où $U=X$. Dans le cas d'un $U$ quelconque, on note que la question est locale sur $S$, donc on peut supposer $S$ affine. Alors tout Module cohérent sur $X$ est quotient d'un Module cohérent \unsure{localemejt|localement} libre, (puisque $\underline O_X(1)$ est un Module inversible absolument ample sur $X$). Comme \unsure{He|le} Module dual $H'=\operatorname{Hom}(H,\underline O_U)$

